\section{The comparisons and their remaining extensions}

The common local factor determines a scalar after the measure or
operator normalization is fixed. The rational Tate construction
assembles these scalars and proves their functional equation. The
equilibrium construction specifies positive states, with a phase
simplex above the critical inverse temperature and uniqueness at
and below it. Its common Toeplitz restriction leaves the polynomial
augmentation distinct from both a thermal state and the vacuum.

For the prime determinants, the continued objects are normalized germs
on explicit Riemann surfaces. The local exponents are total divisor
contributions, and descent is controlled by periods on all closed
curves. On a smaller positive half-plane, several orders have a
connected simultaneous cover determined by the vector of their
period characters. The reciprocal-real poles generate its full image;
no assumption on coincidences or simplicity of nonreal zeta zeros is
needed. The degree-twenty-four surface for orders two and three is
an explicit case where a connected component of a normalized fibre
product carries more information than either coordinate degree.
For higher orders the unbounded reciprocal-prime ramification
obstructs a finite realization on the whole positive half-plane and
near zero. It establishes no natural-boundary theorem along the
entire imaginary axis.

The formal coefficient algebra and the analytic determinant agree
through an actual evaluation homomorphism on a weighted Dirichlet
Banach algebra. The sharp threshold is the same as the Schatten
threshold, while scalar continuation leaves this coefficient domain
unchanged. The row-finite formal operator algebra admits neither a
continuous extension of finite-rank trace nor an algebraic cyclic
extension with its rank-one normalization. These conclusions concern
the specified recipients. They leave open comparison maps using a
different coefficient category, representation or trace ideal.

The integral divisor bar retains multiplication, augmentation,
coalgebra maps and every finite prime-set inclusion. The torsion
Hall construction supplies a separate geometric product and its
normalized adjoints. The arithmetic finite-boundary construction
supplies a degree-three trace, a global torsor point and finite
gluing, with the positive divisor orientation. A bare cone boundary
has the negative connecting sign. These facts distinguish a
cohomological boundary point from local primitives and from a
prepared state.

Three further constructions would change the arithmetic consequences.
First, the map $\tau_S$ of Open Problem~\ref{problem:bar} must satisfy
its Maurer--Cartan equation, converge in the declared filtration and
commute with enlargement of the set of places. This would turn the
integral divisor model into an arithmetic twisted complex. Second,
an independently defined positive representation must recover the
full Weil functional on a dense admissible convolution domain;
the positive narrow-window subspaces do not give that map. Third,
an arithmetic cohomology theory must produce the operators and exact
spectral multiplicities in Theorem~\ref{thm:deninger} and
Corollary~\ref{cor:hp}. Its degree-one determinant would then equal
$\xi_{\mathbb R}(s)/(2\sqrt2\pi^2)$ with the normalization proved here.
Finite arithmetic gluing supplies no continuous gauge measure or
determinant map; those additional state and comparison data are
specified in Open Problem~\ref{problem:acs}.
